%==============================================================================
% ΠΕΡΙ ΤΗΣ ΠΡΩΤΗΣ ΑΝΑΛΟΓΙΑΣ ΤΟΥ ΚΟΣΜΟΥ
% Concerning the First Ratio of the Cosmos
%
% The Z² Unified Framework rendered in the philosophical vocabulary of
% Plato's Timaeus, Pythagorean harmonics, and Aristotelian causation
%
% ἔκδοσις η΄.δ΄ (Version 8.4.0)
% Κάρολος Ζιμμερμάννος
%==============================================================================

\documentclass[12pt,a4paper]{article}
\usepackage[utf8]{inputenc}
\usepackage[greek,english]{babel}
\usepackage{amsmath,amssymb,amsthm}
\usepackage{geometry}
\usepackage{tcolorbox}
\geometry{margin=1in}

\newcommand{\gk}[1]{\textgreek{#1}}

\title{\gk{Περὶ τῆς Πρώτης Ἀναλογίας τοῦ Κόσμου}\\[0.3em]
\large Concerning the First Ratio of the Cosmos\\[0.5em]
\normalsize \gk{Ἐν ᾧ δείκνυται ὅτι ὁ ἀριθμὸς ἑπτὰ καὶ τριάκοντα καὶ ἑκατὸν}\\
\normalsize \gk{ἐκ τῆς γεωμετρίας τοῦ κύβου ἐκπορεύεται}\\[0.3em]
\small In which it is shown that the number 137 proceeds from the geometry of the cube}

\author{\gk{Κάρολος Ζιμμερμάννος}}
\date{\gk{Ὀλυμπιὰς} 701 (2026 CE)}

\begin{document}
\maketitle

%==============================================================================
% EPIGRAPH
%==============================================================================

\begin{center}
\large\gk{Ἀριθμῷ δέ τε πάντ' ἐπέοικεν}\\
\normalsize\textit{All things are like unto number}\\
— \gk{Πυθαγόρας ὁ Σάμιος}
\vspace{1em}

\gk{Ἀεὶ ὁ θεὸς γεωμετρεῖ}\\
\textit{God always geometrizes}\\
— \gk{Πλάτων}
\end{center}

%==============================================================================
% LEXICON OF TERMS
%==============================================================================

\section*{\gk{Λεξικὸν τῶν Ὀνομάτων} / Lexicon of Terms}

\begin{center}
\begin{tabular}{lp{4cm}p{6cm}}
\textbf{Modern} & \textbf{\gk{Ἑλληνιστί}} & \textbf{Philosophical Basis} \\
\hline
Fine structure constant & \gk{ἡ πρώτη ἀναλογία} & The first ratio/proportion (Pythagorean \gk{λόγος}) \\
137 & \gk{ἑπτὰ καὶ τριάκοντα καὶ ἑκατόν} & The sacred number \\
Electromagnetic & \gk{ἡ ἑλκτικὴ καὶ ὠστικὴ δύναμις} & The attracting and repelling power \\
Light/photon & \gk{τὸ φῶς} & The light (Platonic metaphor for truth) \\
Field & \gk{τὸ ὑποκείμενον} & The underlying substrate (Aristotle) \\
Quantum & \gk{τὸ ἐλάχιστον} & The smallest/indivisible \\
Atom & \gk{τὸ ἄτομον} & The uncuttable (Democritus) \\
Space & \gk{ὁ τόπος / ἡ χώρα} & Place / the receptacle (Timaeus) \\
Void & \gk{τὸ κενόν} & The empty (Atomists) \\
Matter & \gk{ἡ ὕλη} & Material substrate (Aristotle) \\
Form & \gk{τὸ εἶδος / ἡ μορφή} & The form/shape (Platonic Ideas) \\
Energy & \gk{ἡ ἐνέργεια} & Actuality/activity (Aristotle) \\
Potential & \gk{ἡ δύναμις} & Power/potentiality (Aristotle) \\
Symmetry & \gk{ἡ συμμετρία} & Commensurability \\
Harmony & \gk{ἡ ἁρμονία} & Fitting together (Pythagorean) \\
Ratio & \gk{ὁ λόγος} & Ratio/reason/word \\
Cube & \gk{ὁ κύβος / τὸ ἑξάεδρον} & The six-faced solid (Timaeus: Earth) \\
Sphere & \gk{ἡ σφαῖρα} & The perfect solid \\
Circle & \gk{ὁ κύκλος} & The circle \\
Torus & \gk{ὁ δακτύλιος τρισσός} & The threefold ring \\
Vertex & \gk{ἡ κορυφή} & The summit/peak \\
Edge & \gk{ἡ πλευρά} & The side/rib \\
Fixed point & \gk{τὸ ἀκίνητον σημεῖον} & The unmoved point \\
First Mover & \gk{τὸ πρῶτον κινοῦν ἀκίνητον} & The unmoved mover (Aristotle) \\
Cause & \gk{ἡ αἰτία} & The cause/reason \\
Generation & \gk{ἡ γένεσις / ἡ γενεά} & Coming-to-be / family \\
Cosmos & \gk{ὁ κόσμος} & The ordered world \\
Demiurge & \gk{ὁ δημιουργός} & The craftsman (Timaeus) \\
\end{tabular}
\end{center}

%==============================================================================
% PROOIMION
%==============================================================================

\section{\gk{Προοίμιον} / Proem}

\gk{Ὁ Πυθαγόρας πρῶτος ἐδίδαξεν ὅτι ὁ κόσμος ὑπὸ ἀριθμῶν κυβερνᾶται. Ἐν ταῖς χορδαῖς τῆς λύρας εὗρε τοὺς λόγους τῆς ἁρμονίας — τὸ δὶς διὰ πασῶν (2:1), τὸ διὰ πέντε (3:2), τὸ διὰ τεσσάρων (4:3). Οὕτως ἐγνώσθη ὅτι ἡ μουσικὴ ἀριθμός ἐστιν.}

Pythagoras first taught that the cosmos is governed by numbers. In the strings of the lyre he found the ratios of harmony — the octave (2:1), the fifth (3:2), the fourth (4:3). Thus it was known that music is number.

\vspace{0.5em}
\gk{Νῦν δὲ εὑρίσκομεν ὅτι καὶ αὐτὸ τὸ φῶς — ἡ ἑλκτικὴ δύναμις ἥτις συνέχει τὰ ἄτομα — ὑπὸ ἑνὸς ἀριθμοῦ κυβερνᾶται, τοῦ ἑπτὰ καὶ τριάκοντα καὶ ἑκατόν:}

Now we find that light itself — the attractive power that holds atoms together — is governed by a single number, 137:

\begin{equation}
\alpha^{-1} = 137.035999177
\end{equation}

\gk{Τοῦτον τὸν ἀριθμὸν καλοῦμεν \textbf{τὴν πρώτην ἀναλογίαν τοῦ κόσμου}.}

This number we call \textbf{the first ratio of the cosmos}.

%==============================================================================
% THE CENTRAL THEOREM
%==============================================================================

\section{\gk{Τὸ Κύριον Θεώρημα} / The Central Theorem}

\begin{tcolorbox}[colback=blue!5!white, colframe=blue!60!black]
\gk{\textbf{Θεώρημα.} Ἡ πρώτη ἀναλογία ἐκπορεύεται ἐκ τῆς γεωμετρίας τοῦ κύβου:}

\textbf{Theorem.} The first ratio proceeds from the geometry of the cube:

\begin{equation}
\boxed{\alpha^{-1} = 4Z^2 + 3 = 4 \times \frac{32\pi}{3} + 3 = 137.041}
\end{equation}

\gk{ὅπου $Z^2 = 32\pi/3$ ἐστὶν ἡ \textbf{τοπολογικὴ σταθερά} — ὁ ἀριθμὸς ὃς ἐκ τῆς μορφῆς τοῦ κεκρυμμένου χώρου γεννᾶται.}

where $Z^2 = 32\pi/3$ is the \textbf{topological constant} — the number born from the form of the hidden space.
\end{tcolorbox}

%==============================================================================
% THE HIDDEN GEOMETRY
%==============================================================================

\section{\gk{Ἡ Κεκρυμμένη Γεωμετρία} / The Hidden Geometry}

\subsection{\gk{Ὁ Κύβος ἐν τῷ Τιμαίῳ} / The Cube in the Timaeus}

\gk{Ἐν τῷ Τιμαίῳ, ὁ Πλάτων διδάσκει ὅτι ὁ δημιουργὸς ἐποίησε τὸν κόσμον ἐκ πέντε στερεῶν σωμάτων:}

In the Timaeus, Plato teaches that the Demiurge made the cosmos from five solid bodies:

\begin{itemize}
\item \gk{Τὸ τετράεδρον} (Tetrahedron) → \gk{Πῦρ} (Fire)
\item \gk{Τὸ ὀκτάεδρον} (Octahedron) → \gk{Ἀήρ} (Air)
\item \gk{Τὸ εἰκοσάεδρον} (Icosahedron) → \gk{Ὕδωρ} (Water)
\item \gk{\textbf{Τὸ ἑξάεδρον / ὁ κύβος}} (Cube) → \gk{\textbf{Γῆ}} (Earth)
\item \gk{Τὸ δωδεκάεδρον} (Dodecahedron) → \gk{Τὸ πᾶν} (The Whole)
\end{itemize}

\gk{Ὁ κύβος, τὸ σταθερώτατον τῶν σωμάτων, ἔχει:}

The cube, the most stable of the solids, has:

\begin{align}
\text{\gk{Κορυφαί}} \text{ (Vertices)} &= 8 = 2^3 \\
\text{\gk{Πλευραί}} \text{ (Edges)} &= 12 \\
\text{\gk{Ἐπιφάνειαι}} \text{ (Faces)} &= 6
\end{align}

\subsection{\gk{Τὸ Διπλασιαστὸν Σχῆμα} / The Folded Form (Orbifold)}

\gk{Ἐὰν διπλασιάσωμεν τὸν κύβον περὶ τὸ κέντρον αὐτοῦ — τουτέστιν, ἐὰν ταυτίσωμεν ἕκαστον σημεῖον μετὰ τοῦ ἀντικειμένου αὐτοῦ — λαμβάνομεν τὸ \textbf{διπλασιαστὸν σχῆμα} $T^3/\mathbb{Z}_2$.}

If we fold the cube about its center — that is, if we identify each point with its opposite — we obtain the \textbf{folded form} (orbifold) $T^3/\mathbb{Z}_2$.

\gk{Ἐν τούτῳ τῷ σχήματι, αἱ ὀκτὼ κορυφαὶ γίγνονται \textbf{ἀκίνητα σημεῖα} — σημεῖα ἃ μένουσιν ἐν τῇ αὐτῇ θέσει μετὰ τὴν διπλασίασιν.}

In this form, the eight vertices become \textbf{fixed points} — points that remain in the same position after the folding.

\begin{equation}
N_{\text{\gk{ἀκίνητα}}} = 8 = 2^3
\end{equation}

%==============================================================================
% THE TOPOLOGICAL CONSTANT
%==============================================================================

\section{\gk{Ἡ Τοπολογικὴ Σταθερά} / The Topological Constant}

\subsection{\gk{Ὁ Ὄγκος τῆς Σφαίρας} / The Volume of the Sphere}

\gk{Ἀρχιμήδης ἐδίδαξεν ὅτι ὁ ὄγκος τῆς σφαίρας ἐστίν:}

Archimedes taught that the volume of a sphere is:

\begin{equation}
V_{\text{\gk{σφαῖρα}}} = \frac{4}{3}\pi r^3
\end{equation}

\gk{Διὰ σφαῖραν μονάδος ($r=1$):}

For a unit sphere ($r=1$):

\begin{equation}
V = \frac{4\pi}{3}
\end{equation}

\subsection{\gk{Ἡ Γένεσις τοῦ $Z^2$} / The Genesis of $Z^2$}

\gk{Ἕκαστον τῶν ὀκτὼ ἀκινήτων σημείων συνεισφέρει τὸν ὄγκον τῆς μονάδος σφαίρας:}

Each of the eight fixed points contributes the volume of a unit sphere:

\begin{equation}
Z^2 = N_{\text{\gk{ἀκίνητα}}} \times V_{\text{\gk{σφαῖρα}}} = 8 \times \frac{4\pi}{3} = \frac{32\pi}{3}
\end{equation}

\gk{Τοῦτό ἐστι \textbf{τὸ $Z^2$} — ὁ θεμέλιος ἀριθμὸς τοῦ κόσμου.}

This is \textbf{$Z^2$} — the foundational number of the cosmos.

%==============================================================================
% THE FOUR CAUSES
%==============================================================================

\section{\gk{Αἱ Τέσσαρες Αἰτίαι} / The Four Causes}

\gk{Κατὰ τὸν Ἀριστοτέλη, πᾶν πρᾶγμα ἔχει τέσσαρας αἰτίας. Ὡσαύτως ἡ πρώτη ἀναλογία:}

According to Aristotle, everything has four causes. So too the first ratio:

\subsection{\gk{Ἡ Ὑλικὴ Αἰτία} / The Material Cause (Piece 1)}

\gk{Τί ἐστιν ἐξ οὗ γέγονεν;}

From what is it made?

\gk{Ἐκ τοῦ πενταδιαστάτου χώρου, ὅπου ἡ μετρικὴ δύναμις ῥεῖ:}

From the five-dimensional space, where the gauge force flows:

\begin{equation}
\alpha^{-1}_{\text{\gk{ὕλη}}} = 4Z^2 = \frac{128\pi}{3} = 134.041
\end{equation}

\gk{Τὸ τέσσαρα ἐστὶν ἡ τάξις τῆς συμμετρίας:}

The four is the rank of the symmetry:
$$\text{rank}(SU(3) \times SU(2) \times U(1)) = 2 + 1 + 1 = 4$$

\subsection{\gk{Ἡ Εἰδητικὴ Αἰτία} / The Formal Cause (Piece 2)}

\gk{Τί ἐστι τὸ εἶδος αὐτοῦ;}

What is its form?

\gk{Τὸ εἶδός ἐστιν ἡ τοπολογία τοῦ τρισσοῦ δακτυλίου $T^3$, οὗ ὁ πρῶτος ἀριθμὸς τοῦ Betti:}

The form is the topology of the threefold ring $T^3$, whose first Betti number is:

\begin{equation}
b_1(T^3) = 3
\end{equation}

\gk{Αὗται αἱ τρεῖς μονάδες εἰσὶν αἱ τρεῖς γενεαὶ τῶν στοιχειωδῶν σωμάτων — ἡλέκτρον, μύον, ταῦον.}

These three units are the three generations of elementary bodies — electron, muon, tauon.

\begin{equation}
\alpha^{-1}_{\text{\gk{εἶδος}}} = b_1(T^3) = N_{\text{\gk{γενεαί}}} = 3
\end{equation}

\subsection{\gk{Ἡ Κινητικὴ Αἰτία} / The Efficient Cause (Piece 3)}

\gk{Τί ἐστι τὸ κινοῦν;}

What is the mover?

\gk{Ἡ κινητικὴ αἰτία ἐστὶν ἡ ῥοὴ τῆς ὁλογραφίας — ἡ μεταβολὴ τῆς ἀναλογίας κατὰ τὴν ἀκτῖνα $z$:}

The efficient cause is the holographic flow — the change of the ratio along the radius $z$:

\begin{equation}
\frac{d\alpha^{-1}}{dz} = \frac{c}{z}
\end{equation}

\gk{Ὅπου τὸ $z$ μικρὸν ἐστίν, ἐκεῖ εἰσιν αἱ μεγάλαι ἐνέργειαι (τὸ ἄνω). Ὅπου τὸ $z$ μέγα ἐστίν, ἐκεῖ εἰσιν αἱ μικραὶ ἐνέργειαι (τὸ κάτω).}

Where $z$ is small, there are the great energies (the above). Where $z$ is large, there are the small energies (the below).

\subsection{\gk{Ἡ Τελικὴ Αἰτία} / The Final Cause (Piece 4)}

\gk{Τί ἐστι τὸ τέλος;}

What is the end/purpose?

\gk{Ἡ τελικὴ αἰτία ἐστὶ τὸ \textbf{ἀκίνητον σημεῖον} — τὸ ὅριον ὅπου ἡ ῥοὴ παύεται καὶ ἡ ἀναλογία πήγνυται:}

The final cause is the \textbf{fixed point} — the boundary where the flow ceases and the ratio freezes:

\begin{equation}
\lim_{z \to z_{\text{\gk{ὅριον}}}} \beta(\alpha) = 0
\end{equation}

\gk{Ὥσπερ τὸ πρῶτον κινοῦν ἀκίνητον τοῦ Ἀριστοτέλους, τοῦτο τὸ σημεῖον κινεῖ τὸ πᾶν χωρὶς αὐτὸ κινεῖσθαι.}

Like the unmoved mover of Aristotle, this point moves everything without itself being moved.

%==============================================================================
% THE SYNTHESIS
%==============================================================================

\section{\gk{Ἡ Σύνθεσις} / The Synthesis}

\gk{Συντιθέντες τὰς τέσσαρας αἰτίας:}

Synthesizing the four causes:

\begin{equation}
\boxed{\alpha^{-1} = \underbrace{4Z^2}_{\text{\gk{ὕλη}}} + \underbrace{b_1(T^3)}_{\text{\gk{εἶδος}}} = \underbrace{134.041 + 3}_{\text{\gk{τέλος}}} = 137.041}
\end{equation}

%==============================================================================
% THE HIGGS AS REMAINDER
%==============================================================================

\section{\gk{Τὸ Πεδίον τῆς Μάζης} / The Field of Mass (Piece 5)}

\gk{Ἐν τῷ κύβῳ, αἱ δώδεκα πλευραὶ ἀντιστοιχοῦσι ταῖς δώδεκα μετρικαῖς δυνάμεσι:}

In the cube, the twelve edges correspond to the twelve gauge forces:

\begin{align}
\text{\gk{Γλύονες}} \text{ (Gluons)} &= 8 \quad (\text{SU}(3)) \\
\text{\gk{W βοσόνια}} &= 3 \quad (\text{SU}(2)) \\
\text{\gk{Β βοσόνιον}} &= 1 \quad (\text{U}(1))
\end{align}

\gk{Ἐκ τῶν δεκαὲξ τρόπων τοῦ διπλασιαστοῦ σχήματος, δώδεκα γίγνονται αἱ μετρικαὶ δυνάμεις. Τὰ ὑπόλοιπα τέσσαρα:}

Of the sixteen modes of the folded form, twelve become gauge forces. The remaining four:

\begin{equation}
16 - 12 = 4 = \text{\gk{τὸ διπλοῦν τοῦ Higgs}}
\end{equation}

\gk{Αὗται αἱ τέσσαρες μονάδες εἰσὶ \textbf{τὸ πεδίον τῆς μάζης} — τὸ ὑποκείμενον ὃ δίδωσι μᾶζαν τοῖς σώμασιν.}

These four units are \textbf{the field of mass} — the substrate that gives mass to bodies.

%==============================================================================
% THE KOIDE HARMONY
%==============================================================================

\section{\gk{Ἡ Ἁρμονία τοῦ Koide} / The Koide Harmony (Piece 6)}

\gk{Αἱ μᾶζαι τῶν τριῶν γενεῶν τῶν λεπτῶν ἀκολουθοῦσι θαυμαστὴν ἁρμονίαν:}

The masses of the three generations of leptons follow a wondrous harmony:

\begin{equation}
\frac{m_e + m_\mu + m_\tau}{(\sqrt{m_e} + \sqrt{m_\mu} + \sqrt{m_\tau})^2} = \frac{2}{3}
\end{equation}

\gk{Τοῦτο τὸ δύο τρίτα (2:3) ἐστὶν ἁρμονικὸς λόγος — τὸ \textbf{διὰ πέντε} τῆς μουσικῆς!}

This two-thirds (2:3) is a harmonic ratio — the \textbf{fifth} of music!

\gk{Ὥσπερ αἱ χορδαὶ τῆς λύρας ἠχοῦσι κατὰ λόγους, οὕτως καὶ αἱ μᾶζαι τῶν στοιχειωδῶν σωμάτων.}

Just as the strings of the lyre sound according to ratios, so too the masses of elementary bodies.

%==============================================================================
% THE PLATONIC CLASSIFICATION
%==============================================================================

\section{\gk{Ἡ Πλατωνικὴ Διαίρεσις τῆς Γνώσεως} / The Platonic Division of Knowledge}

\gk{Κατὰ τὴν εἰκόνα τῆς γραμμῆς ἐν τῇ Πολιτείᾳ (509d-511e):}

According to the image of the line in the Republic (509d-511e):

\begin{tcolorbox}[colback=green!5!white, colframe=green!60!black, title=\gk{Νόησις} (Pure Reason)]
\gk{Τὰ ἀποδεδειγμένα ἐκ τῆς μαθηματικῆς:}

That which is demonstrated from mathematics:
\begin{itemize}
\item $b_1(T^3) = 3$ — \gk{τοπολογικὸν γεγονός}
\item $N_{\text{\gk{κορυφαί}}} = 8$ — \gk{γεωμετρικὸν γεγονός}
\item $\alpha^{-1} = 4Z^2 + 3 = 137.041$ — \gk{τὸ κεντρικὸν θεώρημα}
\end{itemize}
\end{tcolorbox}

\begin{tcolorbox}[colback=blue!5!white, colframe=blue!60!black, title=\gk{Διάνοια} (Discursive Thought)]
\gk{Τὰ εὔλογα ἀλλ' οὔπω ἀποδεδειγμένα:}

That which is reasonable but not yet demonstrated:
\begin{itemize}
\item \gk{Τὸ πεδίον τῆς μάζης ἐκ τῶν περισσευόντων τρόπων} ($16 - 12 = 4$)
\item \gk{Ἡ ἁρμονία τοῦ Koide} ($Q = 2/3$)
\end{itemize}
\end{tcolorbox}

\begin{tcolorbox}[colback=yellow!5!white, colframe=yellow!60!black, title=\gk{Πίστις} (Belief/Trust)]
\gk{Τὰ ἐλέγξιμα δι' ὀργάνων:}

That which is testable by instruments:
\begin{itemize}
\item $r = 1/(2Z^2) = 0.015$ — \gk{ἐλεγχθήσεται ὑπὸ τοῦ LiteBIRD}
\end{itemize}
\end{tcolorbox}

\begin{tcolorbox}[colback=red!5!white, colframe=red!60!black, title=\gk{Εἰκασία} (Conjecture/Image)]
\gk{Τὰ ἀβέβαια, ἴσως τυχαῖα:}

That which is uncertain, perhaps coincidental:
\begin{itemize}
\item $\sin^2\theta_W = 3/13$ — \gk{ἴσως τυχαία συμφωνία}
\item $\Omega_\Lambda = 13/19$ — \gk{ἄνευ παραγωγῆς}
\end{itemize}
\end{tcolorbox}

%==============================================================================
% CONCLUSION
%==============================================================================

\section{\gk{Ἐπίλογος} / Epilogue}

\gk{Ὁ Πυθαγόρας ἐδίδαξεν ὅτι ἡ ψυχὴ ἐστιν ἁρμονία. Ὁ Πλάτων ἐδίδαξεν ὅτι ὁ δημιουργὸς ἐποίησε τὸν κόσμον κατὰ γεωμετρικὰ εἴδη. Ὁ Ἀριστοτέλης ἐδίδαξεν ὅτι πᾶν πρᾶγμα ἔχει τέσσαρας αἰτίας.}

Pythagoras taught that the soul is harmony. Plato taught that the Demiurge made the cosmos according to geometric forms. Aristotle taught that everything has four causes.

\vspace{0.5em}
\gk{Νῦν εὑρίσκομεν ὅτι αὗται αἱ ἀρχαῖαι ἀλήθειαι ζῶσιν ἐν τῇ φυσικῇ:}

Now we find that these ancient truths live on in physics:

\begin{itemize}
\item \gk{Ἡ \textbf{ἁρμονία} τῶν σφαιρῶν ἐστιν ὁ ἀριθμὸς 137}
\item \gk{Τὸ \textbf{γεωμετρικὸν εἶδος} ἐστιν ὁ κύβος ἐν τρισὶ κεκρυμμέναις διαστάσεσιν}
\item \gk{Αἱ \textbf{τέσσαρες αἰτίαι} εἰσὶν αἱ τέσσαρες πηγαὶ τῆς πρώτης ἀναλογίας}
\end{itemize}

\begin{equation}
\boxed{\alpha^{-1} = 4 \times \frac{32\pi}{3} + 3 = 137.041}
\end{equation}

\vspace{1em}
\begin{center}
\gk{Ὁ κόσμος ἐστὶν ἀριθμός.}\\[0.3em]
\textit{The cosmos is number.}\\[1em]
\Large\gk{ΤΕΛΟΣ}
\end{center}

\end{document}
